\section{Related Work}
\label{sec:related-work}

\begin{DeletedBlock}
\paragraph{Stable Domain Theory}

\roly{Now covered in Future Work.} Stable Domain Theory was originally proposed by \citet{berry79} as a
refinement of domain theory aimed at capturing the intensional behaviour of sequential programs, and
elaborated on subsequently by \citet{berry82} and \citet{amadio-curien}. Standard domain-theoretic models
interpret programs as continuous functions, preserving directed joins; Berry observed that this continuity
condition alone is too permissive to model sequentiality. Stability imposes additional constraints to reflect
how functions preserve bounded meets of approximants, effectively requiring that the evaluation of a function
respect a specific computational order. Though stable functions do not fully characterise sequentiality,
because they admit $\mathrm{gustave}$-style counterexamples (\exref{parallel-or}), they remain an appropriate
notion for studying the sensitivity of a program to partial data at a specific point.

Our use of Stable Domain Theory diverges from the traditional aim of modelling infinite or partial data,
however. Instead, we follow a line of work that uses partiality as a qualitative notion of approximation
suitable for provenance and program slicing (discussed in more detail in \secref{related-work:galois-slicing}
below). Paul Taylor’s characterisation of stable functions via local Galois connections on principle downsets
provides the semantic underpinning for the reverse maps used in Galois slicing~\cite{taylor99}. Our work
builds on these ideas by interpreting Galois slicing as a form of differentiable programming, using the
machinery of CHAD to present Galois slicing in a denotational style.
\end{DeletedBlock}

\paragraph{Automatic Differentiation}

Automatic differentiation (AD), discussed in \secref{approx-as-tangents:autodiff}, is the idea of computing
derivatives of functions expressed as programs by systematically applying the chain rule. The observation that
these derivative computations could be interleaved with the evaluation of the original program is due to
\citet{linnainmaa76}, who showed how the forward derivative $\pushf{f}_x$ of $f$ at a point $x$ could be
computed alongside $f(x)$ in a single pass, dramatically improving the efficiency of derivative evaluation
over symbolic or numerical differentiation. This insight became the foundation of forward-mode AD, which
underpins many optimisation and scientific computing tools, including JAX~\cite{jax2018github}.
\citet{griewank89} showed how the Wengert list, the linear record of assignments used in forward-mode to
compute derivatives efficiently, could be traversed in reverse to compute the pullback map. This two-pass
approach is the foundation of reverse-mode AD, and closely resembles implementations of Galois slicing
(\secref{related-work:galois-slicing} below) that record a trace during forward slicing for use in backward
slicing.

More recent approaches to automatic differentiation have emphasised semantic foundations. \citet{elliott18}
proposed a categorical model of AD that interprets programs as functions enriched with their derivatives,
giving a compositional account of differentiation based on duality and linear maps. Vákár and
collaborators~\cite{vakar22,nunes2023} developed the CHAD framework which inspired this paper, using
Grothendieck constructions over indexed categories to capture both values and their tangents in a
compositional semantic structure. These perspectives shed light on the categorical structure of AD and guide
the design of systems that generalise AD% beyond real analysis
, including the application to data provenance
and slicing explored in this paper.

\paragraph{\GPS}
\label{sec:related-work:galois-slicing}

Galois slicing was introduced by Perera and collaborators~\cite{perera12a,perera13} as an operational approach
to program slicing for pure functional programs, based on Galois connections between lattices of input and
output approximations. \Deleted{A connection to stable functions in relation to minimal slices for
short-circuiting operations was alluded to in~\citet{perera13}, but not explored.} Subsequent work extended
the approach to languages with assignment and exceptions~\cite{ricciotti17} and
\Replaced{concurrency}{concurrent systems, applying Galois slicing to the $\pi$-calculus}~\cite{perera16d}.
\Deleted{For the $\pi$-calculus the analysis shifted from functions to transition relations, considering
individual transitions $P \longrightarrow Q$ between configurations $P$ and $Q$ as analogous to the edge
between $x$ and $f(x)$ in the graph of $f$, and building Galois connections between $\downset{P}$ and
$\downset{Q}$.} \Replaced{The present work is instead denotational, recasting dependency analysis as
differentiation over a semiring; the Boolean case recovers Galois slicing, and other semirings give further
analyses. Galois slicing is also more general than our framework in two respects that we plan to consider in
future work (\secref{conclusion}): it computes \emph{program slices} from approximation lattices representing
partially erased programs, and it approximates the \emph{shape} of data rather than only the values at fixed
positions.}{The main difference with the approach presented here is that the earlier work also computes
\emph{program slices}, using approximation lattices that represent partially erased programs; we discuss this
further in \secref{conclusion} below.} More recently, \citet{perera22} presented a variant of Galois slicing
over Boolean algebras\Replaced{, which unlike plain lattices are complemented, giving each analysis a
conjugate}{ rather than plain lattices}. \Replaced{Composing an analysis with its conjugate computes
\emph{related outputs}, which they used to support interactive visualisations with linked selections.
\citet{bond25} revisited this using \emph{dynamic dependence graphs}, computing the conjugate from the
opposite graph.}{In this setting every Galois connection $f \dashv g: A \to B$ has a conjugate $g^\circ \dashv
f^\circ: B \to A$, where $f^\circ$ denotes the De Morgan dual $\neg \comp f \comp \neg$~\cite{jonsson51}. The
provenance analysis can then be composed with its own conjugate to obtain a Galois connection which computes
\emph{related outputs} (e.g., selecting a region of a chart and observing the regions of other charts which
share data dependencies). \citet{bond25} revisited this approach using \emph{dynamic dependence graphs} to
decouple the derivation of dependency information from the analyses that make use of it, and observing that to
compute the conjugate analysis one can just use the opposite graph.} \Added{The related-outputs analysis is
the cognacy relation shown in \secref{examples:linked}, computed here by composing a Boolean Jacobian with its
transpose.}

\paragraph{\added{Semiring Provenance}}

\Added{Provenance for database queries has its own semiring tradition, originating with the \emph{provenance
semirings} of \citet{green07}: a query result carries an annotation from a commutative semiring, and the free
semiring $\mathbbm{N}[X]$ of provenance polynomials is universal, specialising to other semirings by
evaluation. Our framework draws on the same semirings but computes a \emph{derivative} rather than an
annotation: a Jacobian over the semiring relating input positions to output positions, in place of a value
attached to the output. A derivative is a linear map, so the provenance it records is \emph{linearised}
(\secref{approx-as-tangents:linearised-provenance}), and joint use appears as a coefficient rather than the
degree of a polynomial. Over a non-positive semiring, the Boolean dependency our framework reports is a sound
over-approximation, since the exact analysis can cancel a contribution to zero where the Boolean composite
still records a dependency. Extending provenance to deletion and difference likewise calls for non-positive
annotations~\cite{green09,amsterdamer11,geerts10}.}

\paragraph{\added{Discrete and Qualitative Derivatives}}

\Added{Discrete data has its own differential calculi, developed independently of numerical differentiation. The
\emph{Boolean differential calculus} of \citet{steinbach17} arose in the 1950s from the detection of faults in
digital circuits; its \emph{simple derivative} $\partial f / \partial x_i = f|_{x_i = 0} \veebar f|_{x_i = 1}$,
the exclusive-or of $f$ with the input $x_i$ set to $0$ and to $1$, is $1$ exactly when flipping that input
flips the output. Exclusive-or and conjunction make the Booleans a \emph{ring} rather than a lattice, so this is
a derivative over the two-element field, the instance $S = \mathbbm{F}_2$ of our construction. It differs from
the Boolean case of Galois slicing (\secref{related-work:galois-slicing}) only in the semiring. Because
exclusive-or cancels, the derivative is exact, whereas the join of the Boolean lattice only reports a possible
dependency. The sign domain gives a second such calculus. The rule of signs of abstract
interpretation~\cite{cousot77} and the \emph{qualitative derivatives} of qualitative
physics~\cite{dekleer84,kuipers86} ask whether increasing an input increases, decreases, or cannot affect an
output, and are the sign semiring of \secref{approx-as-tangents:qualitative}. Our account treats the semiring as
a parameter, placing these calculi alongside automatic differentiation and provenance as instances of a single
construction.}

\paragraph{Tangent Categories and Differential Linear Logic}

\emph{Tangent categories}, due originally to \citet{rosický84} and developed by \citet{cockett14,cockett18},
provide an abstract categorical framework for reasoning about differentiation, inspired by the structure of
the tangent bundle in differential geometry. In a tangent category, each object $X$ is equipped with a tangent
bundle $T(X)$, and each morphism $f: X \to Y$ has a corresponding differential map $T(f): T(X) \to T(Y)$
satisfying axioms analogous to the chain rule and linearity of differentiation. Tangent categories generalise
Cartesian differential categories \cite{cdcs}, which model differentiation over Cartesian closed categories
using a syntactic derivative operator. Reverse Tangent categories \cite{reverse-tangents} further axiomatise
the existence of reverse derivatives. \Replaced{Our construction realises differentiation of this kind concretely, over an arbitrary commutative
semiring, and we conjecture that the categories underlying it are examples of tangent, or reverse tangent,
categories.}{In \conref{tangent-stable-fns} and
\conref{tangent-fam}, we have conjectured that the categories we have identified in this paper as models of
\GPS are Tangent categories. This would clarify the role of higher derivatives in \GPS, which we conjecture
are related to \emph{program differencing}.} There are likely links to Differential Linear Logic
\cite{ehrhard_differential_2006}. Differential Linear Logic and the Dialectica translation have been used to
model reverse differentiation by \citet{kerjean-pedrot2024}.


% \paragraph{Lens Categories}

% \cite{spivak19}
